\bookchapter{Глава 54. Проект 1 — моя учётная книга расходов (локальная версия)}

\chapterintro{Добро пожаловать в ваш первый «серьёзный проект». Эта учётная книга, хоть и не выходит в сеть, полнокровна: записывает траты, показывает список, сводит статистику по месяцам, умеет удалять и править. Закончив её, вы целиком пройдёте маршрут «данные логика интерфейс». Не спешите срисовывать код — сперва подумайте, потом делайте: это наш вечный принцип.}

\booksection{Думаем о требованиях: сначала задайте пять вопросов}

Прежде чем браться за дело, распутайте замысел с помощью «пяти элементов требований»:

\begin{center}
\begin{tabularx}{\linewidth}{|X|X|}
\hline
\textbf{Элемент} & \textbf{Ответ} \tabularnewline
\hline
Who (для кого) & Для себя; можно и для родных и друзей \tabularnewline
What (что делаем) & Запись доходов и расходов, список операций, сводка за месяц, статистика по категориям \tabularnewline
Why (зачем) & Тренировка + наконец понять, куда уходят деньги \tabularnewline
How (как делаем) & Kotlin + Compose + Room + MVVM, всё локально, без сервера \tabularnewline
Constraint (ограничения) & Уложиться в неделю; минимум — Android 8; даже без публикации писать по стандартам для релиза \tabularnewline
\hline
\end{tabularx}
\end{center}

\begin{center}\small\color{gray}Пять элементов требований учётной книги\end{center}

Дальше разбейте проект на крошечные шаги, каждый из которых «уже полезен»: (1) слой данных Room (можно сохранить запись); (2) экран списка (можно смотреть операции); (3) экран ввода (можно добавить запись); (4) удаление и редактирование; (5) карточка месячной сводки. После каждого шага приложение остаётся «вещью, которой можно пользоваться».

\booksection{Слой данных: даём записям «дом»}

Сначала определите, как выглядит одна запись, — поля нужны в меру, без жадности:

\begin{lstlisting}[language=Kotlin]
@Entity(tableName = "records")
data class Record(
    @PrimaryKey(autoGenerate = true) val id: Long = 0,
    val type: String,        // "income" доход | "expense" расход
    val amount: Double,      // сумма (в юанях)
    val category: String,    // категория: еда, транспорт, зарплата…
    val note: String = "",   // заметка (необязательно)
    val date: String         // "2026-08-15"
)
\end{lstlisting}

Добавьте в Dao два «статистических» запроса — это душа учётной книги:

\begin{lstlisting}[language=Kotlin]
@Dao
interface RecordDao {
    @Query("SELECT * FROM records ORDER BY date DESC, id DESC")
    fun getAll(): Flow<List<Record>>

    @Insert
    suspend fun insert(record: Record)

    @Delete
    suspend fun delete(record: Record)

    @Query("SELECT COALESCE(SUM(amount),0) FROM records WHERE type='income' AND date LIKE :month || '%'")
    suspend fun monthIncome(month: String): Double   // month в виде "2026-08"

    @Query("SELECT COALESCE(SUM(amount),0) FROM records WHERE type='expense' AND date LIKE :month || '%'")
    suspend fun monthExpense(month: String): Double
}
\end{lstlisting}

\begin{tcbinfo}{Совет: почему сумма — Double?}
Строгие финансовые приложения берут BigDecimal или «Int в копейках», чтобы избежать ошибок плавающей точки (старая песня про 0.1+0.2$\neq$0.3). Для тренировочного проекта Double хватает с избытком; если однажды сделаете настоящий финансовый продукт, вспомните про целые «копейки» — об эти грабли споткнулось неисчислимое число предшественников.

\end{tcbinfo}

\booksection{MVVM: интерфейс и логика живут раздельно}

Позволить интерфейсу напрямую трогать базу? Заработает, но код интерфейса разрастётся в зловонную простыню. Архитектура MVVM разводит всех по ролям: View (интерфейс) только рисует, ViewModel (логический штаб) ведёт данными, Model (Room) отвечает за хранение.

View (интерфейс на Compose)

ViewModel (состояние и логика)

Model (слой данных Room)

\begin{center}\small\color{gray}Рисунок 5. Трёхслойная структура MVVM — данные текут в одну сторону, слои друг к другу не лезут\end{center}

\begin{lstlisting}[language=Kotlin]
class RecordViewModel(private val dao: RecordDao) : ViewModel() {

    // Превращаем Flow в наблюдаемое состояние интерфейса
    val records = dao.getAll()
        .stateIn(viewModelScope, SharingStarted.WhileSubscribed(), emptyList())

    // Действия пользователя: вызываем suspend-функции
    fun addRecord(type: String, amount: Double, category: String, note: String, date: String) {
        viewModelScope.launch {
            dao.insert(Record(type = type, amount = amount,
                category = category, note = note, date = date))
        }
    }

    fun deleteRecord(record: Record) {
        viewModelScope.launch { dao.delete(record) }
    }
}
\end{lstlisting}

В интерфейсе (Compose) это работает так: val vm: RecordViewModel = viewModel(), затем val list by vm.records.collectAsState() — данные изменились, список сам обновился, а вам не нужно ничего делать. Вот сила союза MVVM + Flow + Compose.

\booksection{Интерфейс: экран списка + экран ввода}

Для списка берите LazyColumn («ленивый» список) — он не тормозит даже на огромных данных:

\begin{lstlisting}[language=Kotlin]
@Composable
fun RecordListScreen(vm: RecordViewModel) {
    val records by vm.records.collectAsState()

    Column {
        // сводная карточка сверху (доходы - расходы = баланс)
        SummaryCard(records)

        // список операций
        LazyColumn {
            items(records, key = { it.id }) { record ->
                RecordRow(record) { vm.deleteRecord(record) }
            }
        }

        // плавающая кнопка «новая запись»  открывает диалог ввода
        FloatingActionButton(onClick = { showAddDialog = true }) {
            Text("+")
        }
    }
}

@Composable
fun RecordRow(record: Record, onDelete: () -> Unit) {
    Row(Modifier.fillMaxWidth().padding(12.dp)) {
        Text(record.category, fontWeight = FontWeight.Bold)
        Spacer(Modifier.weight(1f))
        Text(if (record.type == "income") "+${record.amount}" else "-${record.amount}",
             color = if (record.type == "income") Color(0xFF27AE60) else Color(0xFFC0392B))
    }
}
\end{lstlisting}

Для ввода проще всего AlertDialog (диалог): несколько полей ввода + переключатель доход/расход + выпадающий список категорий; по «ОК» вызывается vm.addRecord(...).

\begin{tcbwarn}{Внимание: проверяйте ввод, не ленитесь}
Поле суммы проверять обязательно: пустое значение, отрицательное число, «abc» — всё это надо отсекать (удобно через toDoubleOrNull()). Одна из самых частых причин падений приложения — «считать пользовательский ввод примерным отличником». Запомните: всё, что вводит пользователь, проверяйте как врага на границе.

\end{tcbwarn}

\booksection{Редактирование и удаление: данным нужен и «вход», и «выход»}

Учётная книга, которая умеет только добавлять, но не удалять, надоедает за три дня. Добавьте две операции:

\begin{lstlisting}[language=Kotlin]
// Добавим в Dao ещё два метода:
@Query("UPDATE records SET amount=:amount, category=:category, note=:note WHERE id=:id")
suspend fun update(id: Long, amount: Double, category: String, note: String)

@Query("DELETE FROM records WHERE id=:id")
suspend fun deleteById(id: Long)

// В интерфейсе: долгое нажатие на пункт списка открывает меню «правка/удалить» (в Compose — DropdownMenu или AlertDialog)
// Перед удалением лучше показать подтверждение — случайное удаление больше всего ненавидят пользователи финансовых приложений
\end{lstlisting}

Два правила взаимодействия: удаление — только с повторным подтверждением (защита от случайного нажатия), при редактировании по возможности подставляйте прежние значения (исправить один знак проще, чем вводить заново). Не недооценивайте эти мелочи — из таких мелочей и складывается ощущение настоящего продукта.

\booksection{Приёмка: три уровня проверки + два задания со звёздочкой}

\booksection{Приёмка: три уровня проверки + два задания со звёздочкой}

Закончив проект, проведите приёмку по трём уровням «функциональность крайние случаи архитектура»: записывается ли, удаляется ли, остаются ли данные после перезапуска (функциональность); не падает ли при вводе суммы 0 или отрицательного числа, не теряются ли данные при повороте экрана (крайние случаи); не смешались ли интерфейс с данными, нет ли вписанных напрямую значений (архитектура).

Как только базовая версия пройдёт приёмку, добавьте себе два «блюда со звёздочкой»:

\begin{enumerate}
  \item Фильтр по месяцам: сверху переключатель месяца (например, 2026-08 v), и в списке остаются операции только выбранного месяца;
  \item Круговая диаграмма по категориям: через Canvas в Compose или лёгкую библиотеку диаграмм нарисуйте «куда ушли деньги» пирогом — с ней учётная книга уже похожа на продукт.
\end{enumerate}

\begin{tcbpractice}{Практика: доделайте за эту неделю}
Двигайтесь по порядку «слой данных ViewModel экран списка экран ввода карточка сводки», шаг в день. Застряли — швыряйте текст ошибки как есть в поисковик или ИИ (помните седьмой раздел «Практика программирования с ИИ»? Описывая проблему, чётко назовите «ожидаемое поведение, фактическое поведение, шаги воспроизведения»). Когда доделаете, ведите в нём учёт неделю — пользоваться приложением собственного изготовления так приятно, что оно само понесёт вас дальше.

\end{tcbpractice}

\begin{tcbsummary}{Коротко о главном}
\begin{itemize}
  \item Пять элементов требований разбить на маленькие шаги, каждый из которых полезен, — и лишь потом браться за дело
  \item Слой данных: сущность Record + Dao (включая SQL для месячной статистики) + Database
  \item Разделение в MVVM: View только рисует, ViewModel ведёт логику, Model хранит
  \item Flow + stateIn + collectAsState: данные изменились — интерфейс обновился сам
  \item LazyColumn — «ленивый» список; AlertDialog — для ввода
  \item Трёхуровневая приёмка (функциональность / крайние случаи / архитектура) + задания со звёздочкой (фильтр по месяцам, круговая диаграмма)
\end{itemize}

\end{tcbsummary}
